\subsection{Axes IR Definitions and Examples}
\label{app:axes}

\subsubsection{Operation Declarations}
\label{sec:axes-ir}
An Axes IR declaration provides the local mappings used in
Section~\ref{sec:axis-mapping}. Its \code{axes} block binds ordered groups
to operand positions and specifies the output order; its \code{semantics}
block gives the value computed at each output coordinate. In the following
dot declaration, batch and contraction lists pair operand axes in list
order. The remaining axes preserve their operand order.

\begin{lstlisting}[language=TEPL,
  caption={Axes IR representation of generalized dot, with explicit attributes, axis bindings, and value semantics.},
  label={lst:dot-declaration}]
op dot(lhs: tensor, rhs: tensor) -> tensor {
  attrs {
    lb, lc, rb, rc: axis[];
  }
  axes {
    let lf = ordered_complement(all_axes(lhs), attrs.lb ++ attrs.lc);
    let rf = ordered_complement(all_axes(rhs), attrs.rb ++ attrs.rc);
    bind B { lhs: attrs.lb, rhs: attrs.rb };
    bind K { lhs: attrs.lc, rhs: attrs.rc };
    bind I { lhs: lf };
    bind J { rhs: rf };
    yield B ++ I ++ J;
  }
  semantics {
    yield Reduce(K, Product(lhs, rhs));
  }
}
\end{lstlisting}


\subsubsection{Concrete Axis Binding Example}
\label{extended:sec:axis-groups}
Using the attributes in Listing~\ref{lst:dot-declaration}, consider these
operand layouts and bindings:
\begin{lstlisting}
lhs:    [i, k0, b, k1]
rhs:    [k1, j, k0, b]
output: [b, i, j]

lb = [2]
rb = [3]
lc = [1, 3]
rc = [2, 0]
\end{lstlisting}
The batch binding pairs position $2$ of \code{lhs} with position $3$ of
\code{rhs}. The contraction lists pair $(1,2)$ for $k_0$ and $(3,0)$ for
$k_1$. Remaining axes give $\vec I=\langle i\rangle$ and $\vec J=\langle j\rangle$;
$\vec K$ is reduced and therefore omitted from the output. In mapping notation,
\begin{equation}
\begin{aligned}
\mathsf{I}_{\mathrm{lhs}}&=\langle i,k_0,b,k_1\rangle,
 &q_{\mathrm{lhs}}&=\langle1,3\rangle,\\
\mathsf{I}_{\mathrm{rhs}}&=\langle k_1,j,k_0,b\rangle,
 &q_{\mathrm{rhs}}&=\langle2,0\rangle,\\
\vec G_{\mathrm{lhs}}&=\vec G_{\mathrm{rhs}}=\vec K=\langle k_0,k_1\rangle.
\end{aligned}
\end{equation}
Consequently,
\begin{equation}
 \gath(\mathsf{I}_{\mathrm{lhs}},q_{\mathrm{lhs}})
 =\vec K=\gath(\mathsf{I}_{\mathrm{rhs}},q_{\mathrm{rhs}}).
\end{equation}
The resulting computation is
\begin{equation}
 \operatorname{out}[b,i,j]
 =\sum_{k_0,k_1}\operatorname{lhs}[i,k_0,b,k_1]
                    \operatorname{rhs}[k_1,j,k_0,b].
\end{equation}
Paired selections must have equal lengths and corresponding extents,
but equal extents alone do not identify axes. Grouping keeps each axis's
coordinate separate; it does not flatten the axes into one dimension.
Groups may also be empty, as with an unbatched dot's batch group.

\subsubsection{Value and Index Expressions}
\label{extended:sec:value-expressions}
\label{sec:value-expressions}
The following examples and tables define the supported value and index expressions.

Elementwise addition without broadcasting binds every axis of
\code{lhs} to the corresponding axis of \code{rhs}, preserving that group:
\begin{lstlisting}[language=TEPL]
axes {
  bind S {
    lhs: all_axes(lhs),
    rhs: all_axes(rhs)
  };
  yield S;
}
semantics {
  yield Add(lhs, rhs);
}
\end{lstlisting}
The shared binding requires equal rank and matching sizes. At each coordinate
in $\vec S$, the operand references read corresponding elements and \code{Add}
produces their sum. Table~\ref{extended:tab:value-expressions} lists the value primitives.
In this language, \code{Reduce} means summation; it does not imply support
for other reductions such as maximum.

\begin{table}[htbp]
\caption{Value-expression primitives.}
\label{extended:tab:value-expressions}
\begin{tabularx}{\linewidth}{@{}lY@{}}\toprule
Expression & Meaning\\\midrule
\code{input} & Read an operand element using its axis bindings.\\
\code{Const(c)} & Produce a signed integer constant.\\
\code{Add(a, b)} & Add scalar values.\\
\code{Product(a, b)} & Multiply scalar values.\\
\code{Reduce(G, e)} & Sum $e$ over all coordinate combinations of group $\vec G$.\\
\code{Read(input, indices)} & Read an operand at explicitly calculated indices.\\\bottomrule
\end{tabularx}
\end{table}

Index expressions produce an index sequence with one entry for each accessed
axis. Table~\ref{extended:tab:index-expressions} gives their forms. A slice, for example,
uses
\begin{lstlisting}[language=TEPL]
Read(input,
  IndexOffset(attrs.starts,
    IndexScale(attrs.strides, Coords(O))))
\end{lstlisting}
Here $\vec O$ is the output group. Per axis, the index expression computes
\code{starts + strides * output_indices}. The \code{Read} expression retrieves
the selected value. The index calculation determines the location, and the
enclosing value expression determines the scalar result.

\begin{table}[htbp]
\caption{Index-expression primitives.}
\label{extended:tab:index-expressions}
\begin{tabularx}{\linewidth}{@{}lY@{}}\toprule
Expression & Meaning\\\midrule
\code{Coords(G)} & Current coordinates of $\vec G$ in group order.\\
\code{IndexAdd(a, b)} & Add index sequences elementwise.\\
\code{Concat(a, b)} & Concatenate index sequences.\\
\code{IndexScale(attrs.strides, indices)} & Scale each index by its corresponding entry.\\
\code{IndexOffset(attrs.starts, indices)} & Add the corresponding offset to each index.\\\bottomrule
\end{tabularx}
\end{table}


\FloatBarrier
\subsection{Mapping Derivations and Construction Algorithm}
\label{app:mapping}
Symbolic axis mapping connects groups defined by different operations. It
records where each group appears in a tensor, allowing composition of value
expressions while preserving axis order and operand pairings.


\subsubsection{Building and Composing Mapping Tables}
\label{extended:sec:mapping-tables}
This subsection derives the complete axis mapping for
$T=\operatorname{add}(X,Y)$ and $O=\operatorname{dot\_general}(T,Z)$,
including how operand positions preserve axis order. This example composes
addition with dot, extending Section~\ref{sec:axis-mapping}'s dot example.
We assume no broadcasting; ranks and position attributes remain symbolic.
Write $\rho(\vec G)=\operatorname{range}(|\vec G|)$ for all positions in group $\vec G$.

\paragraph{1. Establish addition's shared axis order.}
Addition binds \code{X} and \code{Y} to one ordered group $\vec S$ and yields $\vec S$.
Thus $X,Y,T$ share the same ordered axis interface, with $\vec S$ located at
$\rho(\vec S)$ in each tensor.

\paragraph{2. Locate the dot groups while preserving operand order.}
The dot introduces batch group $\vec B$, contraction group $\vec K$, and free groups
$\vec I,\vec J$. Let $b_\ell,b_r$ be its batch attributes, $k_\ell,k_r$ its contraction
attributes, and $\vec S_Z$ the initial interface of \code{Z}. Free positions are
\begin{equation}
\begin{aligned}
 f_\ell&=\ocomp(\rho(\vec S),b_\ell\catseq k_\ell),\\
 f_r&=\ocomp(\rho(\vec S_Z),b_r\catseq k_r).
\end{aligned}
\end{equation}
Bindings require valid position lists and matching paired extents and must
cover each operand without overlap. Groups can be interleaved, so concatenating
$\vec B,\vec K,\vec I$ does not necessarily reproduce the left operand's dimension order.
Instead, binding positions restore the original order:
\begin{equation}
\begin{aligned}
 \widehat{\vec S}&=\gath(\vec B\catseq \vec K\catseq \vec I,
                  \operatorname{inverse}(b_\ell\catseq k_\ell\catseq f_\ell)),\\
 \widehat{\vec S}_Z&=\gath(\vec B\catseq \vec K\catseq \vec J,
                  \operatorname{inverse}(b_r\catseq k_r\catseq f_r)).
\end{aligned}
\label{extended:eq:rebuilt-interfaces}
\end{equation}
Here $\widehat{\vec S}$ describes the same ordered axes as $\vec S$ using the dot's
groups; it does not transpose the tensor. The same holds for
$\widehat{\vec S}_Z$ and $\vec S_Z$.
For example, if the left operand's axes are $[i,k_0,b,k_1]$, concatenating
$\vec B,\vec K,\vec I$ gives $[b,k_0,k_1,i]$ at positions $[2,1,3,0]$. The inverse
position list restores the original order $[i,k_0,b,k_1]$.
Each concatenated position sequence is a permutation because the bindings
form a complete disjoint partition. Thus
\begin{equation}
\begin{aligned}
 \gath(\widehat{\vec S},b_\ell)&=\vec B,&
 \gath(\widehat{\vec S},k_\ell)&=\vec K,&
 \gath(\widehat{\vec S},f_\ell)&=\vec I.
\end{aligned}
\end{equation}
The right interface satisfies analogous bindings. An existing producer
interface is reused when it already satisfies the bindings.

\paragraph{3. Propagate those positions to \code{X} and \code{Y}.}
Addition describes \code{T} using $\vec S$, whereas the dot describes those same
axes using $\widehat{\vec S}$. Since addition preserves order, each axis of its
output corresponds to the same-position axes in \code{X} and \code{Y}.
Transporting its reads into $\widehat{\vec S}$ therefore places $\vec B,\vec K,\vec I$ at positions
$b_\ell,k_\ell,f_\ell$ in all three tensors. The compiler can now compose
\code{Reduce(K, Product(Add(X,Y), Z))}, with $\vec K$ referring to corresponding
coordinates in \code{X}, \code{Y}, and \code{Z}.

These positions can also be derived through existing mappings. If $q_X$
locates $\widehat{\vec S}$ in $X$, Equation~\eqref{eq:gather-composition} locates
$\vec K$ in $X$ at $\gath(q_X,k_\ell)$. For addition, $q_X$ lists all positions
in order, so the result simplifies to $k_\ell$.

\paragraph{4. Construct the complete mapping table.}
The dot yields $\vec B\catseq \vec I\catseq \vec J$, with $\vec K$ summed out. Combining these
output positions with the operand mappings gives the table below.
Columns denote tensors; the first row gives each axis interface, and the
remaining rows give group positions. A dash indicates no mapping in this
example. The row $\widehat{\vec S}$ represents addition's whole-axis group using
the dot's coordinate description.
\begin{center}
\begin{tabular}{@{}lccccc@{}}\toprule
Interface/group & $X$ & $Y$ & $T$ & $Z$ & $O$\\\midrule
Axis interface & $\widehat{\vec S}$ & $\widehat{\vec S}$ & $\widehat{\vec S}$ & $\widehat{\vec S}_Z$ & $\vec B\catseq \vec I\catseq \vec J$\\
$\widehat{\vec S}$ & $\rho(\widehat{\vec S})$ & $\rho(\widehat{\vec S})$ & $\rho(\widehat{\vec S})$ & -- & --\\
$\vec B$ & $b_\ell$ & $b_\ell$ & $b_\ell$ & $b_r$ & $\rho(\vec B)$\\
$\vec K$ & $k_\ell$ & $k_\ell$ & $k_\ell$ & $k_r$ & --\\
$\vec I$ & $f_\ell$ & $f_\ell$ & $f_\ell$ & -- & $|\vec B|+\rho(\vec I)$\\
$\vec J$ & -- & -- & -- & $f_r$ & $|\vec B|+|\vec I|+\rho(\vec J)$\\\bottomrule
\end{tabular}
\end{center}
Adding a length to a position sequence shifts each entry. These output
offsets therefore follow directly from the declared group order.


\subsubsection{Symbolic Axis Partitioning}
\label{extended:sec:partitioning}
This subsection derives the subgroup positions, operand pairings, and complete
mapping table for the dot-reassociation example in
Section~\ref{sec:partitioning}, extending it to source dots with batch axes.

\paragraph{1. Set up the source mappings.}
Let $T=\operatorname{dot\_general}(X,Y)$ and
$O=\operatorname{dot\_general}(T,Z)$. The inner dot contracts $\vec K$ and
yields $\mathsf{I}_T=\vec B\catseq \vec I\catseq \vec J$, with free groups $\vec I$ from \code{X} and
$\vec J$ from \code{Y}. We assume that the \code{where} guard requires the outer
dot's left batch list to equal the positions of $\vec B$ in $T$, in the same
order. This restricts applicability; the completed rewrite still requires
verification. Write $k_\ell^{\mathrm{outer}},k_r^{\mathrm{outer}}$ for the
outer contraction lists in \code{T} and \code{Z}. Their shared group is
\begin{equation}
 \vec D=\gath(\mathsf{I}_T,k_\ell^{\mathrm{outer}}).
\end{equation}
Since $\vec B$ remains batch, $\vec D$ selects only from $\vec I,\vec J$. In the replacement
\code{dot_general(Y,Z)}, the selected $\vec J$ axes can be contracted, while
the partners of selected $\vec I$ axes must survive for the final dot with \code{X}.

\paragraph{2. Find the selected part of each producer group.}
Write $\rho(\vec G)=\operatorname{range}(|\vec G|)$ for the positions within group $\vec G$.
The producer groups occupy
\begin{equation}
 q_I=|\vec B|+\rho(\vec I),\qquad q_J=|\vec B|+|\vec I|+\rho(\vec J)
\end{equation}
in $\mathsf{I}_T$. Adding an offset to a position list adds it to every entry.
For each $\vec G\in\{\vec I,\vec J\}$, compare its positions with the outer contraction list:
\begin{equation}
 (u_G,v_G)=\mpo(k_\ell^{\mathrm{outer}},q_G).
\end{equation}
For distinct-entry sequences $a,b$, the operation
$\mpo(a,b)=(u,v)$ scans $a$ in order. Whenever $a_i=b_j$, it appends $i$ to
$u$ and $j$ to $v$. It returns positions in the input lists, rather than
the shared values. Hence $\gath(a,u)=\gath(b,v)$, with both results empty
when there is no overlap. Thus $u_G$ records contraction-list positions,
while $v_G$ records positions within $\vec G$.

In the figure, $T$ has axes $[i_0,i_1,j_0,j_1]$ and
$k_\ell^{\mathrm{outer}}=[3,0]$ selects $[j_1,i_0]$.
Matching against $q_I=[0,1]$ gives $u_I=[1]$, $v_I=[0]$;
matching against $q_J=[2,3]$ gives $u_J=[0]$, $v_J=[1]$.
The selected and remaining subgroups are
\begin{equation}
\begin{aligned}
 \vec G_{\mathrm{sel}}&=\gath(\vec G,v_G)=\gath(\vec D,u_G),\\
 w_G&=\ocomp(\rho(\vec G),v_G),\qquad
 \vec G_{\mathrm{rem}}=\gath(\vec G,w_G).
\end{aligned}
\end{equation}
Here $\Isel=[i_0]$, $\Irem=[i_1]$, $\Jsel=[j_1]$, and $\Jrem=[j_0]$.
Selected subgroups follow contraction order; remaining subgroups preserve
producer order. Their concatenation need not restore the original group.
For symbolic inputs, these computations remain expressions over position lists.

\paragraph{3. Map selected axes to \code{Z}.}
The outer dot pairs its left and right lists entry by entry. Selecting slots
$u_I$ and $u_J$ from its right list therefore gives the positions in \code{Z}
paired with the selected subgroups:
\begin{equation}
 z_I=\gath(k_r^{\mathrm{outer}},u_I),\qquad
 z_J=\gath(k_r^{\mathrm{outer}},u_J).
\label{extended:eq:z-partners}
\end{equation}
In the figure, \code{Z} has axes $[z_0,f,z_1]$ and
$k_r^{\mathrm{outer}}=[2,0]$. The positions from step 2 give
$z_I=[0]$ and $z_J=[2]$: $i_0$ pairs with $z_0$, and $j_1$ with $z_1$.
Using contraction-list positions $u_G$, rather than within-group positions
$v_G$, preserves these pairings.

\paragraph{4. Extend the original mapping table.}
Write $b_X,k_X,i_X$ for the known
positions of $\vec B,\vec K,\vec I$ in \code{X}, and $b_Y,k_Y,j_Y$ for those of $\vec B,\vec K,\vec J$ in
\code{Y}. Composition locates each subgroup: for example,
$\gath(i_X,v_I)$ locates $\Isel$ in \code{X}. In \code{Z}, let $b_Z$ be the
batch positions and $f_Z$ the positions of free group $\vec H$.

The inner reduction removes $\vec K$ from $T$; the outer reduction removes
$\Isel,\Jsel$ from $O$. The surviving order is
\begin{equation}
 \mathsf{I}_O=\vec B\catseq\Irem\catseq\Jrem\catseq \vec H.
\label{extended:eq:source-output}
\end{equation}
Table~\ref{extended:tab:refined-mappings} combines the refined relationships. To keep
entries readable, define
\begin{equation}
\begin{aligned}
 i_s&=\gath(i_X,v_I),&i_r&=\gath(i_X,w_I),\\
 j_s&=\gath(j_Y,v_J),&j_r&=\gath(j_Y,w_J),\\
 t_{Is}&=\gath(q_I,v_I),&t_{Ir}&=\gath(q_I,w_I),\\
 t_{Js}&=\gath(q_J,v_J),&t_{Jr}&=\gath(q_J,w_J),\\
 o_I&=|\vec B|+\rho(\Irem),\\
 o_J&=|\vec B|+|\Irem|+\rho(\Jrem),\\
 o_H&=|\vec B|+|\Irem|+|\Jrem|+\rho(\vec H).
\end{aligned}
\end{equation}
These are abbreviations for the position expressions, not additional
attributes to be supplied by the user.
\begin{table}[htbp]
\caption{Refined mappings across the original reassociation graph. Each entry
gives positions in the column's interface. A dash indicates no mapping in
this example.}
\label{extended:tab:refined-mappings}
\centering
\begin{tabular}{@{}lccccc@{}}\toprule
Group/interface & $X$ & $Y$ & $T$ & $Z$ & $O$\\\midrule
Axis interface & $\mathsf{I}_X$ & $\mathsf{I}_Y$ & $\mathsf{I}_T$ & $\mathsf{I}_Z$ & $\mathsf{I}_O$\\
$\vec B$ & $b_X$ & $b_Y$ & $\rho(\vec B)$ & $b_Z$ & $\rho(\vec B)$\\
$\vec K$ & $k_X$ & $k_Y$ & -- & -- & --\\
$\Isel$ & $i_s$ & -- & $t_{Is}$ & $z_I$ & --\\
$\Irem$ & $i_r$ & -- & $t_{Ir}$ & -- & $o_I$\\
$\Jsel$ & -- & $j_s$ & $t_{Js}$ & $z_J$ & --\\
$\Jrem$ & -- & $j_r$ & $t_{Jr}$ & -- & $o_J$\\
$\vec H$ & -- & -- & -- & $f_Z$ & $o_H$\\\bottomrule
\end{tabular}
\end{table}

The procedure generalizes to other producer--consumer pairs: compare position
sequences on the shared interface, identify overlap, and add mappings from
each overlapping portion into both parent groups. Each partition is internally
disjoint, while groups in different partitions can overlap.


\subsubsection{General Construction Algorithm}
\label{extended:sec:mapping-algorithm}
Algorithm~\ref{alg:mapping} expands the procedure outlined in
Section~\ref{sec:mapping-algorithm}. It processes producers before consumers, connects each
operation's bindings to its inputs, and extends mappings to the output.
Overlap introduces the finer subgroup relationships described above.
The input is a source graph and its operation definitions; the output comprises
tensor interfaces, composed value expressions, established mappings,
and partitions.

\begin{paperalgorithm}[H]
\caption{Build interfaces, value expressions, and axis mappings.}
\label{alg:mapping}
\begin{lstlisting}[language=Python]
def build_symbolic_axis_mapping(graph):
    V = initialize_inputs(graph)
    mappings, partitions = set(), {}
    for op in producer_first(graph):
        inputs = [V[t] for t in op.inputs]
        groups, bindings = resolve_bindings(op, inputs)
        aligned = []
        for (axes, value), binding in zip(inputs, bindings):
            aligned_axes = build_interface(axes, binding)
            register_mappings(aligned_axes, binding, mappings)
            refine_and_register(aligned_axes, binding, mappings, partitions)
            value = transport(value, axes, aligned_axes, mappings, partitions)
            aligned.append((aligned_axes, value))
        out_axes = evaluate(op.output_axes, groups)
        binding = output_positions(op, groups)
        register_mappings(out_axes, binding, mappings)
        refine_and_register(out_axes, binding, mappings, partitions)
        V[op] = (out_axes, compose(op.value, aligned, groups))
    return V, mappings, partitions
\end{lstlisting}
\end{paperalgorithm}

\paragraph{State.}
The state is $\Sigma=(V,\mathcal M,\Pi)$. For each tensor $t$,
$V[t]=(\mathsf{I}_t,e_t)$ stores its ordered interface and value expression.
$\mathcal M$ contains individual mapping entries $(\mathsf{I},\vec G,q)$. The partition collection
$\Pi$ records which groups cover an interface completely and without overlap:
\begin{equation}
 \pi=\bigl(\mathsf{I},\langle(\vec G_1,q_1),\ldots,(\vec G_m,q_m)\rangle\bigr).
\end{equation}
An interface can have several partitions because different operations group
its axes differently. A partition changes the grouping used in reasoning;
it does not change tensor shape, rank, or axis order. $\mathcal M$ locates the groups,
whereas $\Pi$ says which groups jointly cover the interface exactly once.

For example, if addition produces $T$ with $\vec S=\langle i,j\rangle$, then
$V[T]=(\mathsf{I}_T,\operatorname{Add}(X,Y))$ with $\mathsf{I}_T=\vec S$,
$\mathcal M$ includes $(\mathsf{I}_T,\vec S,\langle0,1\rangle)$, and $\Pi$ includes the single-member
partition $(\mathsf{I}_T,\langle(\vec S,\langle0,1\rangle)\rangle)$. A consumer may split
$\vec S$ into $\vec I=\langle i\rangle$ and $\vec J=\langle j\rangle$. Then $\mathcal M$ gains their
position mappings and $\Pi$ gains the two-member partition alongside the
original one. A shape $\langle d_i,d_j\rangle$ remains unchanged.

\paragraph{Invariants.}
Every mapping must select its group in the stated order:
\begin{equation}
 (\mathsf{I},\vec G,q)\in \mathcal M\ \Longrightarrow\ \gath(\mathsf{I},q)=\vec G.
 \label{extended:eq:mapping-invariant}
\end{equation}
Every partition member satisfies the same invariant, and
$q_1\catseq\cdots\catseq q_m$ must be a permutation of $\rho(\mathsf{I})$.
Thus each axis occurs exactly once within one partition, although different
partitions may overlap. In the addition example,
$\gath(\langle i,j\rangle,\langle0,1\rangle)=\vec S$, and $\langle0,1\rangle$
covers every position. The sequence $\langle0,0\rangle$ would repeat one
axis and omit the other, violating the partition invariant.



\code{initialize_inputs} gives each graph input an interface and read
expression. \code{producer_first} ensures descriptions are available before
use. \code{resolve_bindings} evaluates symbolic position expressions and
establishes validity and pairing conditions. It reuses producer coordinates
where applicable and otherwise introduces local groups. Each operand's
bindings $(\vec G_i,q_i)$ cover its axes.

\code{register_mappings} adds the individual group-position entries.
\code{build_interface} reuses $\mathsf{I}$ when the required groups already occupy the
required positions. Otherwise it constructs
\begin{equation}
 Q=\gath\bigl(\vec G_1\catseq\cdots\catseq \vec G_m,
        \operatorname{inverse}(q_1\catseq\cdots\catseq q_m)\bigr).
\end{equation}
The inverse arranges collected groups in tensor dimension order.
\code{refine_and_register} registers the new partition and compares it
with earlier partitions on the same interface. It uses \code{matchPositions}
to add overlap mappings. \code{transport} updates the producer's read
coordinates from $\mathsf{I}$ to $Q$ and carries mappings and partitions into that
description for the current consumer. Equal interfaces need no coordinate
update. Finally, \code{evaluate} and \code{output_positions} derive the
output interface and positions from the declared order, and \code{compose}
combines transported expressions using the operation's value definition.

\paragraph{Deriving mappings on demand.}
A missing mapping can be recovered by following established relationships:
\begin{equation}
\begin{aligned}
 (\mathsf{I},\vec G,q)\in \mathcal M\ \land\ (\vec G,\vec H,h)\in \mathcal M
 \quad\Longrightarrow\quad
 \gath(\mathsf{I},\gath(q,h))=\vec H.
\end{aligned}
\end{equation}
Lookup can return $\gath(q,h)$ or register the corresponding entry. Some
relationships propagate during construction and others are derived when
queried; enumerating every possible mapping in advance is unnecessary.
Local bindings and output order establish local mappings, paired position
matches establish overlaps, and Equation~\eqref{eq:gather-composition}
composes them. These operations are designed to maintain the invariant;
they do not dispense with checking the completed rewrite.


\FloatBarrier
\subsection{Attribute Synthesis Derivations and Algorithm}
\label{app:synthesis}
The mapping procedure establishes which logical axes occur in source tensors
and where they occur. Synthesis uses these relationships to fill the
replacement's axis-field holes. It normalizes value expressions, identifies
which axes each replacement operation must preserve or reduce, and queries
mappings to construct its attributes.

Figure~\ref{fig:auto-matching} previews this process for distribution of a dot
over addition. The source is lowered into Axes IR and normalized; the RHS
arithmetic skeleton supplies reads to match before its attributes are known.
Read correspondence allows reuse of source coordinates, after which axis
positions construct the two missing dot descriptors. Here addition has no
broadcasting and preserves paired axis order, so both inferred descriptors
equal the source descriptor. The figure's reads $A[\vec I,\vec K]$, $B[\vec I,\vec K]$, and
$C[\vec K,\vec J]$ abbreviate group coordinates, rather than asserting that groups are
contiguous in physical layouts. The skeleton's question marks denote unresolved
reduction information. Equivalence is checked only after construction.

\paperfigure{1}{figures/Attribute synthesis_ from rules to attributes.pdf}{Matching normalized operand reads supplies coordinate correspondences for constructing replacement attributes. In dot distribution without broadcasting, both replacement dots reuse the source attributes. Group labels denote logical coordinates, and question marks indicate unresolved reductions.}{fig:auto-matching}{Three stages connect a distribution rule and operation semantics to normalized
terms, read matching, position queries, and completed replacement descriptors.}
\FloatBarrier



For attribute reconstruction beyond descriptor reuse, continue the guarded
reassociation example under an arithmetic policy permitting the required laws:
\begin{equation}
 \underbrace{\operatorname{dot}(\operatorname{dot}(X,Y),Z)}_{\mathrm{LHS}}
 \ \longrightarrow\
 \underbrace{\operatorname{dot}(X,\operatorname{dot}(Y,Z))}_{\mathrm{RHS}}.
\end{equation}
The rule supplies the replacement graph. Let $T=\operatorname{dot}(X,Y)$ be
the source intermediate, $U=\operatorname{dot}(Y,Z)$ the replacement
intermediate, and $O,O'$ their respective final results.


\subsubsection{Value Expression Normalization}
\label{extended:sec:normalization}
The graphs combine their reads in different orders. Normalization identifies
corresponding computations before constructing the replacement's attributes.
With output coordinates fixed and indexed reads abbreviated, the LHS is
\begin{equation}
 e_{\mathrm{LHS}}=\sum_{\vec D}\left(\sum_{\vec K} XY\right)Z.
\end{equation}
Each sum ranges over the coordinates of its group. The abbreviated notation
hides, but does not discard, the mappings established in
Section~\ref{sec:axis-mapping}.

TEPL uses a sum-of-products representation similar to SPORES's R-polynomial
form~\cite[Section~2.3, Definition~2.4]{spores}. A term contains a coefficient,
indexed operand reads, and symbolic reduction axes:
\begin{equation}
 t=(c,[r_1,\ldots,r_m],\vec D),\qquad
 \operatorname{value}(t)=\sum_{\vec D} c\prod_{i=1}^{m}r_i.
 \label{extended:eq:normalized-term}
\end{equation}
Each read retains its tensor identity and coordinate expression. $\vec D$ is an
ordered sequence of symbolic axes, and an empty sequence denotes no summation.
The empty product is $1$, permitting constants. A normalized expression is a
list $\mathcal E=[t_1,\ldots,t_n]$ denoting
\begin{equation}
 \operatorname{value}(\mathcal E)=\sum_{j=1}^{n}\operatorname{value}(t_j).
\end{equation}
Reads within a term multiply, separate terms add, and the empty term list
represents zero. Examples, again abbreviating indexed reads, are:
\begin{center}
\begin{tabular}{@{}ll@{}}\toprule
Expression & Normalized term list\\\midrule
$XY$ & \code{[(1, [X, Y], [])]}\\
$X(Y+Z)=XY+XZ$ & \code{[(1, [X, Y], []), (1, [X, Z], [])]}\\
$2XY-Z$ & \code{[(2, [X, Y], []), (-1, [Z], [])]}\\
$\sum_{\vec K} XY$ & \code{[(1, [X, Y], K)]}\\\bottomrule
\end{tabular}
\end{center}
This similarity to SPORES concerns representation; it does not transfer a
completeness result to TEPL's supported reasoning.

Under the permitted multiplication and reduction laws, the example normalizes
to
\begin{equation}
 \sum_{\vec D}\left(\sum_{\vec K} XY\right)Z
 =\sum_{\vec D}\sum_{\vec K} XYZ
 =\sum_{\vec D,\vec K}XYZ.
\end{equation}
The inner reduction contributes reads \code{X}, \code{Y}, and axes $\vec K$.
Multiplication incorporates \code{Z} after checking that it does not depend
on $\vec K$. The outer reduction prepends $\vec D$ after establishing that $\vec D$ and $\vec K$
are disjoint. Thus
\begin{lstlisting}
Value expression:
Reduce(D,
  Product(
    Reduce(K, Product(X, Y)),
    Z))

Normalized term:
t = (1, [X, Y, Z], D ++ K)
\end{lstlisting}

Unknown RHS axis fields mean its actual reduction axes and read-coordinate correspondence
are not yet available. Synthesis instead normalizes its \emph{arithmetic
skeleton}: it expands declared value operations while temporarily omitting
axis information. The skeleton for $X(YZ)$ has coefficient $1$, reads
\code{[X,Y,Z]}, and an empty reduction sequence. That empty sequence is a
property of the skeleton representation, not a claim that the completed RHS
has no reductions.

Synthesis matches each RHS read with the same operand's LHS read and reuses
its coordinates. For several terms, it first matches terms with the same
operands. Multiple possible read-coordinate assignments must be shown to
agree; otherwise synthesis remains unresolved. This correspondence provides
coordinates for construction. Reduction domains, coefficients, and full
value equivalence still have to be checked after attributes are known.


\subsubsection{Determine Which Axes Each Operation Preserves}
\label{extended:sec:keep-axes}


After read matching, the logical axes in each RHS input are known, but their
roles in the replacement may change. The source inner dot contracts $\vec K$;
the source outer dot contracts $\vec D$, partitioned into $\Isel$ and $\Jsel$.
The final output retains $\vec B,\Irem,\Jrem,\vec H$. Source mappings show that
\code{Y} contains $\vec B,\vec K,\Jsel,\Jrem$ and \code{Z} contains
$\vec B,\Isel,\Jsel,\vec H$. These lists specify membership rather than physical order.
To build $U$, the compiler must classify these axes as batch, contraction,
left-free, or right-free. In particular, $\vec K$ must survive $U$ for a later
contraction with \code{X}.

Here $\Axes(x)$ denotes the set of logical axes in tensor $x$'s interface.
For one matched product, $\Axes_{\mathrm{in}}(v)$ and
$\Axes_{\mathrm{out}}(v)$ collect axes used by reads inside and outside
RHS node $v$'s subtree. Within the read-once fragment, the mandatory-axis
definition in Equation~\eqref{eq:mandatory-axes} can be written as
\begin{equation}
 M_v=\Axes_{\mathrm{in}}(v)\cap
 \bigl(\Axes_{\mathrm{out}}(v)\cup\Axes(O)\bigr).
 \label{eq:keep-axes}
\end{equation}
\sidefigure{figures/Infer Axis Roles in U.pdf}{An intermediate preserves axes needed outside its computation or in the final output and reduces the remaining shared axes. The figure applies the mandatory-axis rule to the unbatched replacement intermediate $U$.}{fig:infer-u}{A flow compares axes inside the replacement intermediate with axes needed
outside it and in the output, retaining required axes and contracting the rest.}
This is the same mandatory set $M_v$ used in Section~\ref{sec:canonical-axes};
the live set $P_v$ may additionally include axes that the operation cannot
remove. In the general matching path, requirements for a shared intermediate
must agree across matched products.

Figure~\ref{fig:infer-u} shows the decision for an unbatched instance.
Starting with axes in \code{Y} or \code{Z}, it first preserves those needed
by \code{X}, then those required by the final output. The remaining shared
group, $\Jsel$, can be reduced inside $U$. The figure omits batch group $\vec B$;
the equations below retain it for the guarded general case.

For $U$, the inside set is $\Axes(Y)\cup\Axes(Z)$ and the outside set is
$\Axes(X)$.

Group $\vec K$ occurs inside and outside and so is preserved. $\Jsel$ is needed
only inside and can be contracted. Although $\Irem$ belongs to the final
output, it occurs outside $U$ and reaches the final dot through \code{X}.
At the final dot, $\vec K$ and $\Isel$ are shared by \code{X} and $U$ but absent
from $O$, so they are reduced together.

Using group names as shorthand for all their constituent axes,
\begin{equation}
\begin{aligned}
 \Axes_{\mathrm{in}}(U)
    &=\{\vec B,\vec K,\Isel,\Jsel,\Jrem,\vec H\},\\
 \Axes_{\mathrm{out}}(U)
    &=\{\vec B,\vec K,\Isel,\Irem\},\\
 \Axes(O)&=\{\vec B,\Irem,\Jrem,\vec H\},\\
 M_{U}&=\{\vec B,\vec K,\Isel,\Jrem,\vec H\}.
\end{aligned}
\end{equation}



\paragraph{From required axes to operation groups.}
The mandatory set identifies coordinates that must survive. Assigning coordinates to
operation groups additionally uses which operands contain them and whether
the declared value expression reduces the group. For an RHS node $v$ with
operands $x_1,\ldots,x_m$, let $\Axes_i=\Axes(x_i)$, and let $E_g$ be the
nonempty set of operands bound to declared group $g$. Axes appearing in exactly
those operands form
\begin{equation}
 D_g=\left(\bigcap_{i\in E_g}\Axes_i\right)
     \setminus\left(\bigcup_{i\notin E_g}\Axes_i\right).
\end{equation}
For roles that can reduce or preserve those coordinates, a candidate rule is
\begin{equation}
 \Axes_g=\begin{cases}
 D_g\setminus M_v,&\text{if the value definition reduces }g,\\
 D_g\cap M_v,&\text{otherwise.}
 \end{cases}
 \label{extended:eq:group-assignment}
\end{equation}
This candidate rule is subject to the full operation contract; it is not a
replacement for the canonical live-set rules. For a dot, shared axes in $M_v$
are batch and shared axes outside $M_v$ are contracted. All left-only and
right-only axes remain free, even if they are outside $M_v$. In this example,
the free axes are mandatory too, giving the following assignments:
\begin{center}
\begin{tabular}{@{}lcccc@{}}\toprule
Operation & Batch & Contracted & Left free & Right free\\\midrule
$U=\operatorname{dot}(Y,Z)$ & $\vec B$ & $\Jsel$ & $\vec K,\Jrem$ & $\Isel,\vec H$\\
$O'=\operatorname{dot}(X,U)$ & $\vec B$ & $\vec K,\Isel$ & $\Irem$ & $\Jrem,\vec H$\\\bottomrule
\end{tabular}
\end{center}
Free axes follow operand order. The rule relies on distinguishable declared
group roles; ambiguous assignments remain unresolved.


\subsubsection{Construct Attributes from the Mapping Table}
\label{extended:sec:attribute-construction}
Group membership says which axes a dot uses, but its attributes need ordered
positions in actual operands. A query $\pos(\mathsf{I},\vec G)$ returns $q$ satisfying
$\gath(\mathsf{I},q)=\vec G$, using existing mappings and their compositions. Both sides
of a paired binding query the same \emph{ordered} group, ensuring that
corresponding entries select paired axes.


The following derivation expands the batched program in
Section~\ref{sec:reassociation-program}. The first pair of queries constructs
the inner replacement's contraction lists. Its completed interface $\mathsf{I}_U$ then supplies mappings for
the final dot's queries. Both operands select the ordered contraction group
$\Isel\catseq \vec K$.
The short fields \code{lhs_contract} and \code{rhs_contract} abbreviate the
contraction attributes. The following derivation also includes batch lists.


\paragraph{1. Construct the inner replacement dot.}
For $U=\operatorname{dot}(Y,Z)$, the source table already locates $\vec B$ and
$\Jsel$ in both operands. Its attributes are
\begin{equation}
\begin{aligned}
 b_\ell^U&=b_Y,&b_r^U&=b_Z,\\
 k_\ell^U&=\gath(j_Y,v_J),&k_r^U&=z_J.
\end{aligned}
\end{equation}
By Equation~\eqref{extended:eq:z-partners}, $z_J=\gath(k_r^{\mathrm{outer}},u_J)$,
so both lists select $\Jsel$ in the same order. Free positions are
\begin{equation}
\begin{aligned}
 f_Y^U&=\\
 &\ocomp(\rho(\mathsf{I}_Y),b_Y\catseq k_\ell^U),\\
 f_Z^U&=\\
 &\ocomp(\rho(\mathsf{I}_Z),b_Z\catseq k_r^U),
\end{aligned}
\end{equation}
and the dot's output order constructs
\begin{equation}
\begin{aligned}
 \mathsf{I}_U&=\vec B\catseq\gath(\mathsf{I}_Y,f_Y^U)\\
 &\catseq\gath(\mathsf{I}_Z,f_Z^U).
\end{aligned}
 \label{extended:eq:u-interface}
\end{equation}
The compiler extends mappings to $\mathsf{I}_U$ using this layout. This is necessary
because $\vec K$ and $\Jrem$ can be interleaved in \code{Y}, as can $\Isel$ and
$\vec H$ in \code{Z}. The interface preserves those physical orders instead of
assuming that each surviving group is contiguous.



\paragraph{2. Construct the final replacement dot.}
The final dot contracts $\vec K$ and $\Isel$. Filtering source reduction order
$\vec D\catseq \vec K$ to those axes gives $\Isel\catseq \vec K$: axes from $\vec I$ occur
within $\vec D$, before $\vec K$. Use that same order for both contraction lists:
\begin{equation}
\begin{aligned}
 k_\ell^{O'}&=\gath(i_X,v_I)\catseq k_X,\\
 k_r^{O'}&=\pos(\mathsf{I}_U,\Isel)\catseq\pos(\mathsf{I}_U,\vec K).
\end{aligned}
\end{equation}
A fixed binding determines group order, and inferred paired lists must follow
it. Otherwise the implementation prefers source reduction order for reduced
groups and source output order for retained groups, appending remaining axes
in operand order. Table~\ref{extended:tab:replacement-attributes} completes the two
descriptors.
\begin{table}[htbp]
\caption{Candidate attributes for the replacement dots.}
\label{extended:tab:replacement-attributes}
\begin{tabularx}{\linewidth}{@{}lYY@{}}\toprule
Attribute & $U=\operatorname{dot}(Y,Z)$ & $O'=\operatorname{dot}(X,U)$\\\midrule
Left batch & $b_Y$ & $b_X$\\
Right batch & $b_Z$ & $\rho(\vec B)$\\
Left contraction & $\gath(j_Y,v_J)$ & $\gath(i_X,v_I)\catseq k_X$\\
Right contraction & $z_J$ & $\pos(\mathsf{I}_U,\Isel)$\newline $\catseq\pos(\mathsf{I}_U,\vec K)$\\\bottomrule
\end{tabularx}
\end{table}
Removing the final contraction axes leaves
\begin{equation}
 \mathsf{I}_{O'}=\vec B\catseq\Irem\catseq\Jrem\catseq \vec H=\mathsf{I}_O.
\end{equation}
These attributes therefore have the candidate pairings and output order.
Operation legality and value equivalence still require checking.


\subsubsection{Produce and Verify Attribute Programs}
\label{extended:sec:emit-programs}
The tables contain symbolic group queries, while a rule application must
compute concrete attributes. The compiler resolves each query to expressions
over source attributes and shapes. For the inner replacement's right
contraction list, for example,
\begin{equation}
\begin{aligned}
 u_J&=\mpo\bigl(k_\ell^{\mathrm{outer}},
                   |\vec B|+|\vec I|+\rho(\vec J)\bigr)_1,\\
 k_r^U&=\gath(k_r^{\mathrm{outer}},u_J),
\end{aligned}
\label{extended:eq:source-only-example}
\end{equation}
where subscript $1$ selects the first returned sequence. Group lengths and
ranges are themselves derived from source metadata. This expression constructs
positions at each match without accessing tensor elements. Queries involving
$\mathsf{I}_U$ similarly substitute its layout and compose mappings until only source
metadata remains. Emitted programs contain neither unresolved RHS attributes
nor logical-group queries.

For calculated accesses, the compiler compares supplied coordinate
expressions to establish read correspondence. Synthesis constructs only
axis-position expressions; extent and coordinate fields retain their
supplied expressions or schema defaults. Verification checks that these
expressions are defined and satisfy access bounds and output shape.

\paragraph{Check the completed replacement.}
With attributes filled, the RHS has actual coordinates and reduction domains.
In the example, $U$ reduces $\Jsel$ and the final dot reduces
$\Isel\catseq \vec K$:
\begin{equation}
\begin{aligned}
 e_{\mathrm{RHS}}
 &=\sum_{\Isel,\vec K}X\left(\sum_{\Jsel}YZ\right)\\
 &=\sum_{\Isel,\vec K,\Jsel}XYZ.
\end{aligned}
\end{equation}
The selected subgroups partition $\vec D$, and the mappings preserve coordinate
pairing. The normalized RHS can therefore be compared to
$\sum_{\vec D,\vec K}XYZ$ under the permitted reduction laws. Verification also checks
output order, shape, dtype, and operation legality.

After construction, a fresh verification context reconstructs the completed
rule using the emitted programs. Required conditions must follow from legal
source metadata and the declared guard; construction obligations do not
become new assumptions. On success, the result is a collection of checked
source-dependent programs, compiled once and evaluated on each matched
graph's metadata. Unestablished construction or verification relationships
leave the rule unresolved.
